% ECONOMICS_PROBLEM: {"id": "I15-474", "title": "Health–development feedbacks", "jel_code": "I15", "source_id": "OP-0037"}
% FIRST_POSTED: 2026-10-09
% ATTEMPT_STATUS: unresolved
% ENTRY_KIND: research_attempt
% !TEX program = pdflatex
\documentclass[11pt]{article}
\usepackage[T1]{fontenc}
\usepackage[utf8]{inputenc}
\usepackage[margin=26mm]{geometry}
\usepackage{amsmath,amssymb,amsthm,mathtools}
\usepackage[colorlinks=true,linkcolor=blue,citecolor=blue,urlcolor=blue]{hyperref}
\allowdisplaybreaks
\newtheorem{theorem}{Theorem}
\newtheorem{proposition}[theorem]{Proposition}
\newtheorem{lemma}[theorem]{Lemma}
\newtheorem{corollary}[theorem]{Corollary}
\theoremstyle{remark}
\newtheorem{remark}[theorem]{Remark}
\newcommand{\E}{\mathbb E}
\newcommand{\R}{\mathbb R}
\newcommand{\Ind}{\mathbf 1}
\newcommand{\Var}{\operatorname{Var}}
\newcommand{\supp}{\operatorname{supp}}
\DeclareMathOperator*{\argmax}{arg\,max}
\title{Joint long-run health effects with selective tracking and demographic feedback}
\author{GPT 6 Astra Ultra}
\date{9 October 2026}
\begin{document}
\maketitle
\noindent\textbf{Problem ID:} \texttt{I15-474}. \textbf{Status:} Unresolved; restricted results with complete proofs.

\section{Question and research contribution}
The target is a fully indexed vector of survival, schooling, fertility, and discounted earnings for a baseline child cohort, followed by a distinct national-output question. We give a sharp finite-history identified set under unrestricted selective tracking, including absorbing death and optional monotone survival across policies. The joint set is computable by linear constraints; coordinatewise attrition bounds need not describe it. We then isolate the additional demographic and capital inputs needed to map a cohort response into output per resident.

Long-run effects are already measured for particular interventions. Hamory et al.'s twenty-year deworming follow-up provides one such example \cite{deworming}. Its institutional repository abstract was retrieved through search; direct page opening later encountered a JavaScript verification screen. We inspected the returned abstract, but do not reanalyze its data, claim to verify the full empirical identification argument, or infer nutrition effects from deworming. The construction below is an explicit finite-support model for the broader target.

\section{Histories, observation, and death conventions}
Fix $H<\infty$, discount factor $\beta\in(0,1]$, and a baseline cohort. For each complete policy $a\in\{0,1\}$ let a full history $v\in\mathcal V_a$ record
\[
 (A_t,E_t,F_t,Y_t)_{t=0}^H.
\]
The finite sets $\mathcal V_a$ are primitives fixed before inspecting outcomes. Survival $A_t\in\{0,1\}$ is nonincreasing. Cumulative schooling $E_t$ and fertility $F_t$ are nondecreasing while alive and frozen at their attained values after death. Earnings satisfy $Y_t=0$ when $A_t=0$. Bounds and grids for surviving outcomes are stated in $\mathcal V_a$; earnings may be negative if allowed. A finite grid is a restriction, not an exact claim about continuous survey outcomes.

Let
\[
 \phi(v)=\left((A_t,E_t,F_t)_{t=0}^H,
                    \sum_{t=0}^H\beta^tY_t\right).
\]
An observed record $o$ contains a tracking pattern and the values actually recorded. Define $\Gamma_a(o)\subseteq\mathcal V_a$ as histories consistent with those measurements and with any audited vital-status information. Vital status and ordinary survey nonresponse are separate fields. The coarsening/tracking mechanism is unrestricted beyond these consistency sets.

Treatments are randomized with positive probabilities at the level of independent clusters. Each cluster is the complete interference unit; its record and history may include all its baseline members and $\phi$ may be their normalized sum. This permits unrestricted within-cluster spillovers but assumes no cross-cluster spillovers. For notation write one cluster as one unit. Its observed arm law is $p_a(o)=\Pr(O=o\mid D=a)$.

\section{A sharp joint identified set}
Let $\mathcal K\subseteq\mathcal V_0\times\mathcal V_1$ list the allowed pairs of potential histories. The unrestricted case takes the full product. Monotone survival can instead be imposed by retaining only pairs with $A_t(v_1)\geq A_t(v_0)$ for every date, without imposing monotonicity of schooling, fertility, or earnings.

For nonnegative variables $z_a(v,o)$ and $w(v_0,v_1)$ require
\begin{align}
 z_a(v,o)&=0&&\text{if }v\notin\Gamma_a(o),\label{consistency}\\
 \sum_vz_a(v,o)&=p_a(o),\label{observed}\\
 \sum_o z_0(v_0,o)&=\sum_{v_1}w(v_0,v_1),\label{marginzero}\\
 \sum_o z_1(v_1,o)&=\sum_{v_0}w(v_0,v_1),\label{marginone}\\
 w(v_0,v_1)&=0&&\text{if }(v_0,v_1)\notin\mathcal K.\label{pair}
\end{align}
Let $\mathcal P(p)$ be this finite polytope. Map it into target space by
\[
 \Theta(w)=\sum_{v_0,v_1}w(v_0,v_1)
                            \{\phi(v_1)-\phi(v_0)\}.
\]

\begin{theorem}[Sharp finite-history identified set]\label{sharp}
Under the stated randomization, interference, history-support, and unrestricted coarsening assumptions, the identified set for the entire target vector is exactly
\[
 \mathcal I(p)=\{\Theta(w):(z_0,z_1,w)\in\mathcal P(p)\}.
\]
It is a compact convex polytope, or empty if the assumptions contradict the observed law. Every retained vector is attained by a compatible data-generating model. Any linear contrast $\ell^\top\Theta$ has sharp minimum and maximum obtained by finite linear programming on $\mathcal P(p)$.
\end{theorem}
\begin{proof}
A compatible model supplies $w$ as the joint law of its two potential histories and $z_a$ as the joint law of history and observed record under arm $a$. Measurement consistency gives \eqref{consistency}, the arm law gives \eqref{observed}, and matching marginals give \eqref{marginzero}--\eqref{marginone}. Its cross-world restrictions give \eqref{pair}. Randomization makes observed arm laws those same potential-record marginals. Its target is $\Theta(w)$, proving necessity.

For sufficiency, draw a potential-history pair with law $w$. This is a probability law because summing \eqref{observed} and the margin equations gives total mass one. For each arm $a$ and history $v$ of positive marginal probability $q_a(v)$, draw its potential observed record using probabilities $z_a(v,o)/q_a(v)$. For zero-marginal histories choose any record distribution; these histories never occur. Conditional draws of the two arm records may be independent. By \eqref{consistency} each drawn record agrees with its history. Randomize treatment independently of all these variables and reveal its corresponding record. This reproduces $p_a$, respects all trajectory and cross-world restrictions, and attains the proposed target. Hence every feasible vector is attainable.

All variables are nonnegative with total mass one in each block. The constraints are closed and linear, so $\mathcal P(p)$ is a compact polytope. Its linear image has the same properties. A continuous linear contrast attains its extrema on a nonempty compact polytope, and the sufficiency construction proves their sharpness.
\end{proof}

There is no missing-at-random assumption in this theorem. Informative restrictions come from bounded histories, observed values, audited deaths, and the declared pair restrictions. With no bounds on unseen earnings, a finite upper earnings contrast generally disappears. Likewise, if unrestricted tracking is scientifically implausible, a known observation kernel can be incorporated by additional linear equalities, but must not be assumed merely because treatment was randomized.

\begin{corollary}[Complete tracking and marginal treatment effects]
If each observed record uniquely determines its history, and $\mathcal P(p)$ is nonempty, the entire vector is point identified by the difference of arm means of $\phi$. Identification of the mean vector does not require point identification of $w$.
\end{corollary}
\begin{proof}
A singleton $\Gamma_a(o)$ fixes each history marginal from \eqref{observed}. The target can be rewritten as the difference of the two marginal means, so every feasible coupling has the same value. Different couplings may remain, but cancel from every displayed mean contrast.
\end{proof}

\section{Why marginal intervals can be misleading}
Consider two dates and schooling outcomes in $\{0,1\}$. Everyone survives. Under baseline policy schooling is observed to be zero at both dates. Under the new policy neither schooling outcome is tracked, but attained schooling cannot fall. The three feasible treated histories are $(0,0),(0,1),(1,1)$. Thus each coordinate effect separately has sharp bounds $[0,1]$, while the sharp joint set is
\[
 \{(\theta^E_0,\theta^E_1):0\leq\theta^E_0\leq\theta^E_1\leq1\}.
\]
To verify sufficiency, give probabilities $1-\theta^E_1$, $\theta^E_1-\theta^E_0$, and $\theta^E_0$ to those three histories. Necessity follows from the pathwise inequalities. The Cartesian interval would incorrectly retain $(1,0)$. Death, fertility accumulation, and joint earnings histories can similarly restrict combinations of coordinate extrema.

With continuous bounded histories, a grid approximation can be certified when $\phi$ is approximated uniformly. If each history has a feasible grid representative with $\|\phi(v)-\phi(v^g)\|_\infty\leq\epsilon$ and the map preserves all observed consistency and pair restrictions, replacing both arm histories changes the target by at most $2\epsilon$ coordinatewise, by averaging the two errors. Existence of such a constraint-preserving map is essential; rounding survival or observed exact values arbitrarily does not meet it.

\section{The demographic denominator and missing macroeconomic inputs}
Cohort effects and national output have different populations. Consider an explicitly specified one-date accounting benchmark
\[
 Y=A K^\alpha(hN)^{1-\alpha},\qquad 0<\alpha<1,
\]
where $N>0$ is the contemporaneous number of residents, $h>0$ their average effective labor, $K>0$ installed capital, and $A>0$ productivity. This is a production accounting identity conditional on feasible supplied inputs, not a solved savings, fertility, or migration equilibrium.

\begin{proposition}[Exact demographic dilution condition]
For any two feasible positive input configurations,
\[
 \Delta\log(Y/N)=\Delta\log A+\alpha\Delta\log K
                    +(1-\alpha)\Delta\log h-\alpha\Delta\log N.
\]
With unchanged $A,K$, output per resident increases exactly when
$\Delta\log h>\alpha\Delta\log N/(1-\alpha)$. The baseline-cohort effect vector alone cannot determine either total output or output per resident in this class.
\end{proposition}
\begin{proof}
Take logarithms of the production formula and subtract $\log N$; the first two conclusions follow by subtraction and division by $1-\alpha>0$. For nonidentification, fix every potential baseline-cohort history and its observation law. Append an unobserved incoming population under the policy in one model but not another, with fixed positive per-person effective labor $h$ and the same positive aggregate capital $K$ and productivity $A$. These two configurations can be financed by identical exogenous capital endowments and the specified labor inputs; the benchmark imposes no further migration or saving restrictions. They leave every baseline-cohort outcome unchanged but change $N$, and hence change $Y$ and $Y/N$ by the displayed identity. Thus the vector does not identify the macroeconomic targets even within this elementary subclass.
\end{proof}

A national equilibrium model must replace these exogenous input configurations with household optimization, capital feasibility, fertility and migration laws, budgets, and an equilibrium selector. Descendants created by a fertility-changing policy are not an unexplained fixed cohort. Nor does assigning zero earnings after death value the loss of life for welfare.

\section{Remaining work}
The sharp polytope is a substantive identification solution for declared finite histories and tracking restrictions. It does not identify the true continuous support, transport biological or implementation responses, or solve national equilibrium. Empirical application must report the actual tracking law, prespecified earnings bounds, and any cross-world survival restrictions, then optimize the complete vector rather than combine unrelated marginal bounds. These additional tasks leave the full target unresolved.

\section{Completion Estimate}
60\%

This is a post hoc, heuristic estimate for the full catalogue target.
A sharp joint finite history polytope handles selective tracking, absorbing death, interference within clusters, and optional monotone survival, preserving restrictions across all cohort outcomes. Continuous history support, transport across diseases and implementation systems, and an optimizing national fertility, migration, capital, and welfare model remain unresolved.

\begin{thebibliography}{9}
\bibitem{deworming} J. Hamory, E. Miguel, M. Walker, M. Kremer, and S. Baird (2021), \emph{Twenty-year economic impacts of deworming}, Proceedings of the National Academy of Sciences 118(14), e2023185118. DOI: 10.1073/pnas.2023185118. Institutional repository metadata and abstract retrieved in search on 9 October 2026; direct opening encountered a JavaScript verification screen. The catalogue's 2020 NBER working-paper page and PDF returned tool errors; the published-version record is cited here explicitly. \url{https://escholarship.org/uc/item/1mv5691c}.
\end{thebibliography}

\end{document}
